\documentclass[a4paper,11pt]{article}

\usepackage[T1]{fontenc}
\usepackage[utf8]{inputenc}
\usepackage[french]{babel}
\usepackage{geometry}
\usepackage{booktabs}
\usepackage{array}
\usepackage{tabularx}

\geometry{margin=2.2cm}
\pagestyle{empty}

\title{Alanya tailles à prévoir\\pour l'auto-hébergement des tuiles}
\author{}
\date{}

\begin{document}
\maketitle
\thispagestyle{empty}

\noindent
Voie retenue (raster auto-hébergé) : rendu à la demande + cache nginx.
Les tailles ci-dessous sont des ordres de grandeur pour le dimensionnement disque.

\vspace{1.2em}

{\small
\renewcommand{\arraystretch}{1.35}
\begin{tabularx}{\textwidth}{@{}>{\raggedright\arraybackslash}p{2.8cm}cccc>{\raggedright\arraybackslash}X@{}}
\toprule
\textbf{Zone}
  & \textbf{Données OSM} (\texttt{.pbf})
  & \textbf{Base PostGIS}
  & \textbf{Cache tuiles}
  & \textbf{Total réaliste} \\
\midrule
Cameroun
  & $\sim$0,22\,Go
  & $\sim$2--5\,Go
  & 5--20\,Go
  & $\sim$10--25\,Go \\
Afrique francophone
  & $\sim$2,7\,Go
  & $\sim$15--40\,Go
  & 20--80\,Go
  & $\sim$40--120\,Go \\
Afrique
  & $\sim$7,4\,Go
  & $\sim$40--100\,Go
  & 50--150\,Go
  & $\sim$100--250\,Go \\
Monde
  & $\sim$80--90\,Go
  & centaines de Go à $\sim$1\,To+
  & très élevé
  & souvent 1\,To+ \\
\bottomrule
\end{tabularx}
}

\vspace{1.5em}

\noindent\textit{Notes.}
\begin{itemize}
  \itemsep0.35em
  \item Le fichier \texttt{.pbf} sert surtout à l'import ; en production restent PostGIS et le cache.
  \item Le cache grossit avec l'usage (villes et zooms fréquentés), pas avec toute la carte d'un coup.
  \item Planisphère sans tout héberger : niveaux 0--8 monde $\approx$ 1,7\,Go de tuiles, détail seulement sur le Cameroun.
  \item Cible Alanya aujourd'hui : \textbf{Cameroun}, environ 10--25\,Go de disque.
\end{itemize}

\end{document}
